%\usepackage{makecell} % добавить в преамбулу

\section{Рабочий проект}

\subsection{Спецификация компонентов и классов программ}

\subsubsection{Модули и объекты интерфейса пользователя}

Пользовательский интерфейс программного продукта предназначен для управления процессом сжатия изображений в формате JPEG, задания параметров качества, выбора входных и выходных данных, а также отображения результатов сжатия и вспомогательной информации.

Интерфейс включает в себя следующие основные компоненты:
\begin{itemize}
	\item главное окно приложения с элементами управления процессом сжатия;
	\item область выбора входного изображения или каталога изображений;
	\item элементы задания параметра качества JPEG;
	\item область отображения информации о размере файлов до и после сжатия;
	\item индикатор выполнения процесса;
	\item окно визуализации графика зависимости размера файла от качества.
\end{itemize}

Программный интерфейс реализован с использованием библиотеки \textbf{Tkinter} языка программирования Python и ориентирован на интуитивно понятное взаимодействие с пользователем.

\subsubsection{Описание объектов интерфейса программы}

На основе требований, сформулированных в техническом задании, с использованием языка программирования Python и библиотеки Tkinter разработан графический пользовательский интерфейс программы сжатия JPEG-изображений.

На рисунке~\ref{fig:ui_main} представлен интерфейс взаимодействия пользователя с главным окном приложения.

%\clearpage
%\begin{figure}[h]
%	\centering
%	\includegraphics[width=0.9\textwidth]{ui_main.jpg}
%	\caption{Интерфейс взаимодействия пользователя с главным окном приложения}
%	\label{fig:ui_main}
%\end{figure}

В таблице~\ref{tab:ui_objects} приведено описание объектов интерфейса.

\begin{xltabular}{\textwidth}{|c|l|l|X|}
	\caption{Описание объектов интерфейса пользователя\label{tab:ui_objects}}\\
	\hline
	 № & Тип объекта & Имя объекта & Назначение объекта \\
	\hline
	\endfirsthead
	
	\hline
	 № & Тип объекта & Имя объекта & Назначение объекта \\
	\hline
	\endhead
	
	\hline
	\endfoot
	
	1 & Button & select\_file\_btn & Выбор одного JPEG-файла для сжатия \\
	\hline
	2 & Button & select\_folder\_btn & Выбор каталога с JPEG-изображениями \\
	\hline
	3 & Label & input\_display & Отображение пути к выбранному файлу или папке \\
	\hline
	4 & Button & \makecell[tl]{low\_quality\_btn \\ medium\_quality\_btn \\ high\_quality\_btn \\ custom\_quality\_btn} & Кнопки выбора пресета качества \\
	\hline
	5 & Scale & quality\_slider & Ручная настройка качества JPEG в диапазоне от 1 до 100~\% \\
	\hline
	6 & Button & select\_output\_btn & Выбор каталога сохранения сжатых изображений \\
	\hline
	7 & Label & output\_display & Отображение пути выбранного каталога сохранения \\
	\hline
	8 & Label & \makecell[tl]{size\_before\_label \\ size\_after\_label} & Отображение размера файла до и после сжатия \\
	\hline
	9 & ProgressBar & progress & Индикатор выполнения процесса сжатия \\
	\hline
	10 & Button & compress\_btn & Запуск процесса сжатия изображения или набора изображений \\
	\hline
	11 & Button & graph\_btn & Построение графика зависимости размера файла от параметра качества JPEG \\
	\hline
\end{xltabular}

\renewcommand{\arraystretch}{1.0}

\subsubsection{Функциональная спецификация программных модулей}

\paragraph{Модуль encoder}

Модуль реализует основной алгоритм JPEG-сжатия.

\begin{xltabular}{\textwidth}{|X|p{5cm}|p{7cm}|}
	\caption{Функции модуля encoder\label{tab:encoder_funcs}}\\
	\hline
	\centrow Название функции & Аргументы & Назначение \\
	\hline
	\endfirsthead
	
	\hline
	\centrow Название функции & Аргументы & Назначение \\
	\hline
	\endhead
	
	\hline
	\endfoot
	
	encode\_jpeg & image, quality & Основная функция JPEG-кодирования изображения с заданным параметром качества \\
	\hline
	split\_into\_blocks & channel & Разбиение двумерного массива яркости или цветоразностного канала на блоки размером $8\times8$ \\
	\hline
	encode\_channel & channel, qtable & Применение дискретного косинусного преобразования и квантования ко всем блокам канала \\
	\hline
	decode\_channel & blocks, qtable, shape & Обратное восстановление канала изображения из квантованных блоков \\
	\hline
\end{xltabular}

\renewcommand{\arraystretch}{1.0}

\newpage
\paragraph{Модуль пользовательского интерфейса (ui)}

\begin{xltabular}{\textwidth}{|X|p{5cm}|p{7cm}|}
	\caption{Функции модуля ui\label{tab:ui_funcs}}\\
	\hline
	\centrow Название функции & Аргументы & Назначение \\
	\hline
	\endfirsthead
	
	\hline
	\centrow Название функции & Аргументы & Назначение \\
	\hline
	\endhead
	
	\hline
	\endfoot
	
	create\_ui & root & Инициализация графического пользовательского интерфейса и создание элементов управления \\
	\hline
	update\_progress & value & Обновление значения индикатора выполнения процесса сжатия \\
	\hline
	update\_sizes & before, after & Отображение размеров исходного и сжатого файлов \\
	\hline
	show\_error & message & Вывод диалогового окна с сообщением об ошибке \\
	\hline
\end{xltabular}

\renewcommand{\arraystretch}{1.0}

\paragraph{Модуль сохранения (save\_jpeg)}

\begin{xltabular}{\textwidth}{|X|p{5cm}|p{7cm}|}
	\caption{Функции модуля save\_jpeg\label{tab:save_funcs}}\\
	\hline
	\centrow Название функции & Аргументы & Назначение \\
	\hline
	\endfirsthead
	
	\hline
	\centrow Название функции & Аргументы & Назначение \\
	\hline
	\endhead
	
	\hline
	\endfoot
	
	save\_image & image, path & Сохранение восстановленного изображения в формате JPEG по заданному пути \\
	\hline
	prepare\_image & image & Приведение изображения к формату и параметрам, совместимым с сохранением в JPEG \\
	\hline
	calc\_file\_size & path & Определение размера сохранённого файла на диске \\
	\hline
\end{xltabular}

\renewcommand{\arraystretch}{1.0}

\paragraph{Вспомогательные модули}

\begin{xltabular}{\textwidth}{|X|X|p{8cm}|}
	\caption{Функции вспомогательных модулей\label{tab:helper_funcs}}\\
	\hline
	\centrow Модуль & \centrow Функция & Назначение \\
	\hline
	\endfirsthead
	
	\hline
	\centrow Модуль & \centrow Функция & Назначение \\
	\hline
	\endhead
	
	\hline
	\endfoot
	
	color & rgb\_to\_ycbcr & Преобразование изображения из цветового пространства RGB в пространство YCbCr \\
	\hline
	color & ycbcr\_to\_rgb & Обратное преобразование изображения из YCbCr в RGB \\
	\hline
	dct & dct2 & Реализация прямого двумерного дискретного косинусного преобразования \\
	\hline
	dct & idct2 & Реализация обратного двумерного дискретного косинусного преобразования \\
	\hline
	quantization & scale\_quant\_table & Масштабирование стандартных таблиц квантования в зависимости от качества JPEG \\
	\hline
	subsampling & subsample & Выполнение субдискретизации цветоразностных каналов изображения \\
	\hline
	subsampling & upsample & Восстановление исходного разрешения цветоразностных каналов (апсемплинг) \\
	\hline
\end{xltabular}

\renewcommand{\arraystretch}{1.0}

\subsection{Тестирование программной системы}

\subsubsection{Позитивные тестовые сценарии}

\paragraph{Случай 1 — пользователь открыл программу}

Предусловие: программный продукт установлен, все программные зависимости удовлетворены.

Тестовый случай: пользователь запускает программу сжатия JPEG-изображений из командной строки или с помощью ярлыка операционной системы.

Ожидаемый результат: программа отображает главное окно приложения с элементами управления процессом сжатия и пустыми областями отображения входных и выходных данных.

Результат представлен на рисунке~\ref{fig:test1}.

%\begin{figure}[H]
%	\centering
%	\includegraphics[width=0.9\textwidth]{test1}
%	\caption{Запуск программы}
%	\label{fig:test1}
%\end{figure}

\paragraph{Случай 2 — сжатие одного изображения}

Предусловие: программа запущена, выбран корректный JPEG-файл и каталог сохранения.

Тестовый случай: пользователь выбирает один JPEG-файл, задаёт параметр качества и запускает процесс сжатия.

Ожидаемый результат: выбранный файл успешно сжимается, в интерфейсе отображаются размеры файла до и после сжатия.

Результат представлен на рисунке~\ref{fig:test2}.

%\begin{figure}[H]
%	\centering
%	\includegraphics[width=0.9\textwidth]{test2}
%	\caption{Сжатие одного изображения}
%	\label{fig:test2}
%\end{figure}

\paragraph{Случай 3 — сжатие набора изображений}

Предусловие: программа запущена, выбран каталог с JPEG-изображениями и каталог сохранения.

Тестовый случай: пользователь выбирает папку с изображениями и запускает процесс сжатия.

Ожидаемый результат: все изображения из выбранного каталога успешно обработаны, индикатор выполнения корректно отображает ход процесса.

Результат представлен на рисунке~\ref{fig:test3}.

%\begin{figure}[H]
%	\centering
%	\includegraphics[width=0.9\textwidth]{test3}
%	\caption{Сжатие набора изображений}
%	\label{fig:test3}
%\end{figure}

\paragraph{Случай 4 — построение графика зависимости размера от качества}

Предусловие: выбрано одно JPEG-изображение.

Тестовый случай: пользователь нажимает кнопку построения графика зависимости размера файла от параметра качества.

Ожидаемый результат: отображается график зависимости размера JPEG-файла от заданного значения качества сжатия.

Результат представлен на рисунке~\ref{fig:test4}.

%\begin{figure}[H]
%	\centering
%	\includegraphics[width=0.9\textwidth]{test4}
%	\caption{Построение графика зависимости размера от качества}
%	\label{fig:test4}
%\end{figure}

\subsubsection{Негативные тестовые сценарии}

\paragraph{Случай 5 — запуск сжатия без выбора файла}

Предусловие: программа запущена, входной файл или каталог не выбран.

Тестовый случай: пользователь нажимает кнопку запуска сжатия без выбора входных данных.

Ожидаемый результат: отображается сообщение об ошибке с уведомлением о необходимости выбора файла или каталога.

Результат представлен на рисунке~\ref{fig:test5}.

%\begin{figure}[H]
%	\centering
%	\includegraphics[width=0.9\textwidth]{test5}
%	\caption{Запуск сжатия без выбора файла}
%	\label{fig:test5}
%\end{figure}

\paragraph{Случай 6 — не выбран каталог сохранения}

Предусловие: выбран входной файл или каталог, но не указан путь сохранения.

Тестовый случай: пользователь запускает процесс сжатия без выбора каталога сохранения.

Ожидаемый результат: процесс сжатия не запускается, пользователю выводится предупреждающее сообщение.

Результат представлен на рисунке~\ref{fig:test6}.

%\begin{figure}[H]
%	\centering
%	\includegraphics[width=0.9\textwidth]{test6}
%	\caption{Не выбран каталог сохранения}
%	\label{fig:test6}
%\end{figure}

\paragraph{Случай 7 — построение графика без выбора изображения}

Предусловие: входное изображение не выбрано.

Тестовый случай: пользователь нажимает кнопку построения графика.

Ожидаемый результат: выводится предупреждающее сообщение о необходимости выбора изображения.

Результат представлен на рисунке~\ref{fig:test7}.

%\begin{figure}[H]
%	\centering
%	\includegraphics[width=0.9\textwidth]{test7}
%	\label{fig:test7}
%\end{figure}